\chapter*{前言和致谢}


这是一本面向操作系统课程用的草稿教材。它通过研究一个名为 xv6 的内核来阐述操作系统的主要概念。Xv6 以 Dennis Ritchie 和 Ken Thompson 的 Unix Version 6 (v6) 为模型。Xv6 大致遵循 v6 的结构和风格，但使用 ANSI C 为 RISC-V 多核处理器实现。

本文应与 xv6 源代码一起阅读，这种方法受 John Lions 的《UNIX 第6版注释（Lions' Commentary on UNIX 6th Edition）》启发；文本有指向源代码的超链接，位于 \url{https://github.com/mit-pdos/xv6-riscv}。参见 \url{https://pdos.csail.mit.edu/6.1810} 以获取有关 v6 和 xv6 的额外在线资源，包括使用 xv6 的几个实验作业。

我们在麻省理工学院的操作系统课程 6.828 和 6.1810 中使用了本文。我们感谢这些课程的教师、助教和学生，他们都以直接或间接的方式为 xv6 做出了贡献。特别是，我们要感谢 Adam Belay、Austin Clements 和 Nickolai Zeldovich。最后，我们要感谢那些通过邮件向我们报告文本中的错误或提出改进建议的人：Abutalib Aghayev、Sebastian Boehm、brandb97、Anton Burtsev、Raphael Carvalho、Tej Chajed、Brendan Davidson、Rasit Eskicioglu、Color Fuzzy、Wojciech Gac、Giuseppe、Tao Guo、Haibo Hao、Naoki Hayama、Chris Henderson、Robert Hilderman、Eden Hochbaum、Wolfgang Keller、Paweł Kraszewski、Henry Laih、Jin Li、Austin Liew、lyazj@github.com、Pavan Maddamsetti、Jacek Masiulaniec、Michael McConville、m3hm00d、Mes0903、miguelgvieira、Mark Morrissey、Muhammed Mourad、Harry Pan、Harry Porter、Siyuan Qian、Zhefeng Qiao、Askar Safin、Salman Shah、Huang Sha、Vikram Shenoy、Adeodato Simó、Ruslan Savchenko、Pawel Szczurko、Warren Toomey、tyfkda、tzerbib、Vanush Vaswani、Chen Wang、Xi Wang、Zou Chang Wei、Sam Whitlock、Qiongsi Wu、LucyShawYang、ykf1114@gmail.com 和 Meng Zhou。

如果您发现错误或有改进建议，请发送邮件给 Frans Kaashoek 和 Robert Morris（kaashoek,rtm@csail.mit.edu）。
